\documentclass[11pt,letterpaper]{article}
% This source is intentionally restricted to pdfLaTeX.
\usepackage[T1]{fontenc}
\usepackage[utf8]{inputenc}
\usepackage{textcomp}
\usepackage{lmodern}
\usepackage[margin=0.82in,top=0.94in,bottom=0.91in,headheight=18pt]{geometry}
\usepackage{amsmath,amssymb}
\usepackage[table]{xcolor}
\definecolor{DeepNavy}{HTML}{14283F}
\definecolor{Teal}{HTML}{0A747C}
\definecolor{Gold}{HTML}{C78F36}
\definecolor{Slate}{HTML}{536477}
\definecolor{Pale}{HTML}{F4F7FA}
\definecolor{CodeEdge}{HTML}{D4DFE9}
\usepackage{microtype}
\usepackage{parskip}
\setlength{\parskip}{0.47em}
\usepackage{titlesec}
\titleformat{\section}{\normalfont\Large\bfseries\color{DeepNavy}}{\thesection}{0.62em}{}
\titlespacing*{\section}{0pt}{2.0ex plus .7ex minus .2ex}{0.75ex}
\titleformat{\subsection}{\normalfont\large\bfseries\color{Teal}}{\thesubsection}{0.6em}{}
\usepackage{longtable,booktabs,array,calc}
\renewcommand{\arraystretch}{1.22}
\AtBeginEnvironment{longtable}{\small}
\usepackage{etoolbox}
\makeatletter
\patchcmd\longtable{\par}{\if@noskipsec\mbox{}\fi\par}{}{}
\makeatother
\usepackage{listings}
\IfFileExists{upquote.sty}{\usepackage{upquote}}{}
\lstdefinestyle{suitecode}{
  basicstyle=\ttfamily\footnotesize,
  breaklines=true,breakatwhitespace=false,
  columns=fullflexible,keepspaces=true,showstringspaces=false,
  frame=single,rulecolor=\color{CodeEdge},
  backgroundcolor=\color{Pale},
  keywordstyle=\color{Teal}\bfseries,
  commentstyle=\color{Slate}\itshape,
  stringstyle=\color{DeepNavy},
  xleftmargin=5pt,xrightmargin=5pt,
  framexleftmargin=5pt,framexrightmargin=5pt,
  aboveskip=0.85em,belowskip=0.95em,
  tabsize=2
}
\lstset{style=suitecode}
\usepackage{enumitem}
\setlist[itemize]{leftmargin=1.6em,itemsep=2pt,topsep=3pt}
\setlist[enumerate]{leftmargin=1.8em,itemsep=2pt,topsep=3pt}
\providecommand{\tightlist}{\setlength{\itemsep}{0pt}\setlength{\parskip}{0pt}}
\newcommand{\passthrough}[1]{#1}
\usepackage{xurl}
\usepackage[colorlinks=true,linkcolor=Teal,urlcolor=Teal,citecolor=Teal,pdfborder={0 0 0}]{hyperref}
\urlstyle{same}
\usepackage{fancyhdr,lastpage}
\pagestyle{fancy}
\fancyhf{}
\fancyhead[L]{\footnotesize\sffamily\color{Slate} ECMA-262:2025 \textbullet{} Engineering Documentation}
\fancyhead[R]{\footnotesize\sffamily\color{Slate} 16th Edition}
\fancyfoot[L]{\footnotesize\sffamily\color{Slate} Independent technical companion}
\fancyfoot[R]{\footnotesize\sffamily\color{Slate}\thepage\ / \pageref{LastPage}}
\renewcommand{\headrulewidth}{0.35pt}
\renewcommand{\footrulewidth}{0pt}
\setlength{\emergencystretch}{3em}
\setcounter{tocdepth}{2}
\setcounter{secnumdepth}{2}
\begin{document}
\begin{titlepage}
\thispagestyle{empty}
\vspace*{2.8cm}
{\color{Gold}\rule{18mm}{3pt}}\par\vspace{18mm}
{\sffamily\Large\bfseries\color{Teal} SOFTWARE ENGINEERING DOCUMENTATION SUITE}\par
\vspace{8mm}
{\sffamily\huge\bfseries\color{DeepNavy} ECMAScript 2025 }\par
\vspace{4mm}
{\sffamily\fontsize{24}{29}\selectfont\bfseries\color{DeepNavy} Engineering Curriculum }\par
\vspace{10mm}
{\large\color{Slate} Twelve-week prerequisite-based learning path }\par
\vspace{16mm}
{\color{Teal}\rule{\textwidth}{0.8pt}}\par
\vfill
{\sffamily\color{Slate}\large ECMA-262:2025 \textbullet{} 16th Edition \textbullet{} June 2025}\par
\vspace{3mm}
{\sffamily\color{Slate} October 2026 \textbullet{} Document 8 of 8}\par
\vspace{3mm}
{\small\color{Slate} Original engineering companion. The uploaded ECMA-262:2025 specification remains the authoritative source.}\par
\end{titlepage}
\pagenumbering{arabic}
\tableofcontents
\clearpage
\section{Curriculum charter}\label{curriculum-charter}

\textbf{Audience:} software engineers, JavaScript platform engineers, security reviewers, technical educators, and language implementers. \textbf{Reference:} uploaded ECMA-262:2025, 16th edition, June 2025. \textbf{Duration:} 12 weeks, 6--9 hours per week (recommended). \textbf{Outcomes:} trace ECMAScript semantics to normative clauses, write reliable JavaScript using 2025 language features, design tests for edge cases, distinguish host interfaces, and perform language-aware security reviews.

\textbf{Scope classification:} \textbf{{[}SPEC{]}} is standard-defined; \textbf{{[}HOST{]}} is platform-provided; \textbf{{[}EDU{]}} is instructional design or engineering guidance. This curriculum is not itself a certification or conformance program established by Ecma International.

\section{Prerequisite graph and pacing}\label{prerequisite-graph-and-pacing}

\begin{lstlisting}
W01 Specification literacy
  -> W02 Values and conversion
  -> W03 Source, syntax, and binding
  -> W04 References, objects, and classes
  -> W05 Functions, closures, and generators
  -> W06 Scripts and modules
  -> W07 Iterators and collections
  -> W08 Promises and async execution
  -> W09 Built-ins, binary data, and regex
  -> W10 Security review
  -> W11 Conformance testing and CI
  -> W12 Capstone defense
\end{lstlisting}

\begin{longtable}[]{@{}
  >{\raggedright\arraybackslash}p{(\columnwidth - 6\tabcolsep) * \real{0.2500}}
  >{\raggedright\arraybackslash}p{(\columnwidth - 6\tabcolsep) * \real{0.2500}}
  >{\raggedright\arraybackslash}p{(\columnwidth - 6\tabcolsep) * \real{0.2500}}
  >{\raggedright\arraybackslash}p{(\columnwidth - 6\tabcolsep) * \real{0.2500}}@{}}
\toprule\noalign{}
\begin{minipage}[b]{\linewidth}\raggedright
Module
\end{minipage} & \begin{minipage}[b]{\linewidth}\raggedright
Required prior knowledge
\end{minipage} & \begin{minipage}[b]{\linewidth}\raggedright
Standard clauses
\end{minipage} & \begin{minipage}[b]{\linewidth}\raggedright
Deliverable
\end{minipage} \\
\midrule\noalign{}
\endhead
\bottomrule\noalign{}
\endlastfoot
W01 & programming basics & 1--5, 17, Annexes & specification reading worksheet \\
W02 & W01 & 6--7, 21 & coercion trace notebook \\
W03 & W01--W02 & 8, 11--14 & parser/binding test pack \\
W04 & W02--W03 & 9--10, 15, 20, 28 & object semantics demonstration \\
W05 & W04 & 15, 27 & closure/generator lab \\
W06 & W03--W05 & 16, Annex D & module graph and host analysis \\
W07 & W02, W05 & 23--24, 27.1 & lazy iterator/set algebra service \\
W08 & W05--W07 & 9, 27.2, 27.7 & async failure model \\
W09 & W02, W07 & 21--22, 23, 25 & regex/binary fidelity tests \\
W10 & W04--W09 & mechanism clauses, external CWE & secure code audit \\
W11 & W01--W10 & clause 2, relevant API clauses & test traceability matrix \\
W12 & all & integrated & production-style capstone \\
\end{longtable}

\section{Week 1 --- specification literacy}\label{week-1-specification-literacy}

\textbf{Learning objectives:} identify scope, conformance rules, normative-optional and legacy markings, normative/informative annexes, and algorithm conventions. Read clauses 1--5 and Annex B overview. \textbf{Practice:} mark five statements as standard requirement, host decision, or personal style rule; explain \passthrough{\lstinline!?!} and \passthrough{\lstinline"!"} in abstract algorithms. \textbf{Evidence:} one-page concept map. \textbf{Exit criteria:} accurately classify 8/10 provided scenarios and locate their source clauses.

\section{Week 2 --- types and conversions}\label{week-2-types-and-conversions}

\textbf{Objectives:} distinguish eight language types and specification types; trace \passthrough{\lstinline!ToPrimitive!}, \passthrough{\lstinline!ToString!}, \passthrough{\lstinline!ToNumeric!}, equality algorithms, NaN, signed zero, and BigInt. \textbf{Practice:} build a 20-case table for \passthrough{\lstinline!===!}, \passthrough{\lstinline!==!}, \passthrough{\lstinline!Object.is!}, \passthrough{\lstinline!Map!} key equality. \textbf{Lab:} reject invalid monetary input and demonstrate unsafe Number conversion versus BigInt. \textbf{Exit criteria:} complete three annotated conversion traces without unsupported shortcuts.

\section{Week 3 --- parsing and declarations}\label{week-3-parsing-and-declarations}

\textbf{Objectives:} explain lexical goals, ASI, restricted productions, early errors, scope analysis, and lexical initialization. \textbf{Practice:} classify a collection of short snippets as valid script, valid module, early error, or runtime failure. \textbf{Lab:} ASI and temporal-dead-zone probes. \textbf{Exit criteria:} identify all parse-time versus runtime errors in a 12-case test battery.

\section{Week 4 --- object invariants and classes}\label{week-4-object-invariants-and-classes}

\textbf{Objectives:} explain ordinary internal methods, property descriptors, prototype lookup, proxies, and private names. \textbf{Practice:} trace \passthrough{\lstinline![[Get]]!} / \passthrough{\lstinline![[Set]]!} across object prototypes. \textbf{Lab:} implement a private-state class and a proxy that logs access without violating invariants. \textbf{Exit criteria:} explain a failed nonconfigurable proxy trap and create an executable regression test.

\section{Week 5 --- functional semantics}\label{week-5-functional-semantics}

\textbf{Objectives:} distinguish \passthrough{\lstinline!this!} for arrows vs ordinary functions, closures, class methods, generators, async generators, and \passthrough{\lstinline!finally!} semantics. \textbf{Practice:} implement a lazy generator and observe cleanup when \passthrough{\lstinline!for...of!} breaks. \textbf{Lab:} bounded ID generator with early termination. \textbf{Exit criteria:} draw a state-machine trace for \passthrough{\lstinline!next!}, \passthrough{\lstinline!return!}, and \passthrough{\lstinline!throw!}.

\section{Week 6 --- modules and host interfaces}\label{week-6-modules-and-host-interfaces}

\textbf{Objectives:} explain \passthrough{\lstinline!Module Record!}, live bindings, linking, evaluation, top-level \passthrough{\lstinline!await!}, and import attributes. \textbf{Practice:} draw a dependency graph with circular imports and a JSON import. \textbf{Lab:} make the same module run in a target host; document the resolution strategy, supported attributes, and expected failure behavior. \textbf{Exit criteria:} label every host-specific step; never attribute file resolution solely to ECMA-262.

\section{Week 7 --- Iterator helpers and Set algebra}\label{week-7-iterator-helpers-and-set-algebra}

\textbf{Objectives:} work with \passthrough{\lstinline!Iterator.from!}, lazy helpers, terminal consumption, upstream close, and newly standardized Set operations. \textbf{Practice:} build an audit stream using \passthrough{\lstinline!filter!}, \passthrough{\lstinline!map!}, \passthrough{\lstinline!take!}, and \passthrough{\lstinline!toArray!}. \textbf{Lab:} implement permission-set difference and intersection with edge cases. \textbf{Exit criteria:} prove tested behavior for empty inputs, duplicate elements, and early termination. \textbf{{[}SPEC{]}} Iterator and Set 2025 additions are defined in clauses 27.1 and 24.2.

\section{Week 8 --- promises and job scheduling}\label{week-8-promises-and-job-scheduling}

\textbf{Objectives:} distinguish synchronous execution, asynchronous job reactions, thenable assimilation, rejection, and host event-loop integration. \textbf{Practice:} trace a log-ordering example with \passthrough{\lstinline!.then!} and \passthrough{\lstinline!await!}. \textbf{Lab:} normalize sync and async functions via \passthrough{\lstinline!Promise.try!}, then test thrown and rejected outcomes. \textbf{Exit criteria:} identify which ordering claims are language-level and which depend on the host.

\section{Week 9 --- text, numbers, and binary structures}\label{week-9-text-numbers-and-binary-structures}

\textbf{Objectives:} use \passthrough{\lstinline!RegExp.escape!}, reason about inline modifiers, Unicode strings, typed arrays, DataView, and Float16 precision. \textbf{Practice:} build dynamic literal matching without regex injection, test unpaired surrogates, and quantify half-precision rounding. \textbf{Lab:} encode/decode binary metrics using explicit endianness. \textbf{Exit criteria:} demonstrate at least ten edge-case tests and explain each expected result.

\section{Week 10 --- language-aware security assessment}\label{week-10-language-aware-security-assessment}

\textbf{Objectives:} link coercion, getters, prototype mutation, eval sinks, regex compilation, and async race windows to concrete threat models. \textbf{{[}EDU{]}} CWE and OWASP provide external vulnerability vocabulary; not every risky language operation is an exploitable vulnerability. \textbf{Lab:} complete a security assessment of a deliberately flawed request processor; document dataflow, impact, fix, and regression test. \textbf{Exit criteria:} each claimed vulnerability has a reproducible path and test.

\section{Week 11 --- conformance and CI}\label{week-11-conformance-and-ci}

\textbf{Objectives:} distinguish language conformance, host compatibility, and application security testing. \textbf{Practice:} inspect Test262's repository structure and interpretation rules. \textbf{Lab:} add a reproducible CI workflow with unit tests, linting, and a 2025 feature-support matrix; draft a clause-to-test traceability report. \textbf{Exit criteria:} no skipped feature is silently represented as passing.

\section{Week 12 --- capstone engineering defense}\label{week-12-capstone-engineering-defense}

\textbf{Problem:} build an ESM-based policy-evaluation library consuming validated JSON records, filtering data through iterators, computing set relations, handling sync/async rules with consistent promises, and exposing testable error semantics. Deploy it in a documented target host; record support for relevant ECMAScript 2025 APIs.

\textbf{Required deliverables:} (1) threat model and public API contract; (2) source code with no untrusted dynamic evaluation; (3) 30 or more unit/property/negative tests; (4) clause-to-behavior traceability; (5) CI workflow and runtime support matrix; (6) five-minute architecture presentation with trade-offs; (7) remediation notes for at least three deliberately injected defects.

\section{Hands-on problem bank}\label{hands-on-problem-bank}

\begin{longtable}[]{@{}
  >{\raggedright\arraybackslash}p{(\columnwidth - 6\tabcolsep) * \real{0.2500}}
  >{\raggedright\arraybackslash}p{(\columnwidth - 6\tabcolsep) * \real{0.2500}}
  >{\raggedright\arraybackslash}p{(\columnwidth - 6\tabcolsep) * \real{0.2500}}
  >{\raggedright\arraybackslash}p{(\columnwidth - 6\tabcolsep) * \real{0.2500}}@{}}
\toprule\noalign{}
\begin{minipage}[b]{\linewidth}\raggedright
ID
\end{minipage} & \begin{minipage}[b]{\linewidth}\raggedright
Prerequisite
\end{minipage} & \begin{minipage}[b]{\linewidth}\raggedright
Task
\end{minipage} & \begin{minipage}[b]{\linewidth}\raggedright
Passing evidence
\end{minipage} \\
\midrule\noalign{}
\endhead
\bottomrule\noalign{}
\endlastfoot
EX-01 & W02 & distinguish SameValueZero and strict equality & six counterexamples \\
EX-02 & W03 & classify ASI / early errors & parse and runtime matrix \\
EX-03 & W04 & proxy getter invariant & throws expected TypeError \\
EX-04 & W05 & close generator on early break & \passthrough{\lstinline!finally!} observable \\
EX-05 & W06 & module cycle / live export binding & integration test plus graph \\
EX-06 & W07 & compare Set operations & algebraic property tests \\
EX-07 & W07 & lazy iterator bounded consume & upstream \passthrough{\lstinline!.return()!} evidence \\
EX-08 & W08 & normalize callback with Promise.try & value + thrown cases \\
EX-09 & W09 & RegExp.escape literal handling & metachar/Unicode suite \\
EX-10 & W09 & Float16 roundtrip & rounding proof cases \\
EX-11 & W10 & prototype pollution prevention & exploit/fix regression \\
EX-12 & W11 & conformance case traceability & exact source clause and oracle \\
\end{longtable}

\section{Scoring and advancement}\label{scoring-and-advancement}

\begin{longtable}[]{@{}lrl@{}}
\toprule\noalign{}
Component & Weight & Minimum performance \\
\midrule\noalign{}
\endhead
\bottomrule\noalign{}
\endlastfoot
Weekly knowledge checks & 20\% & average 80\% \\
Lab implementations & 30\% & each runnable and reviewed \\
Specification traceability & 20\% & correct clause/algorithm paths \\
Security assessment & 10\% & exploit evidence + remediation \\
Capstone & 20\% & 80\% and no critical design omissions \\
\end{longtable}

\textbf{{[}EDU{]}} These thresholds are instructor-designed, not Ecma requirements. A student should not advance to work involving shared memory, module linking, or language-level security review until prerequisite blocks are demonstrated. Provide remediation assignments for missing concepts rather than granting credit for unsupported explanations.

\section{Instructor delivery plan}\label{instructor-delivery-plan}

For each week: a 60-minute specification walkthrough, 90-minute guided code lab, 60-minute independent exercise, 30-minute peer review, and approximately 2--4 hours independent study. Use explicit rubric feedback. Keep engine compatibility records current; provide alternative tasks if a deployed runtime lacks a 2025 feature, without redefining the normative specification outcome.

\textbf{Rubric dimensions:} correctness, evidence traceability, boundary handling, code clarity, secure behavior, and operational reproducibility. Use observable tests instead of confidence in a particular engine implementation. Ask learners to contrast native behavior with any polyfill rather than claiming a polyfill proves full conformance.

\section{Curriculum governance and versioning}\label{curriculum-governance-and-versioning}

\textbf{{[}EDU{]}} Maintain a source-edition lock (ECMA-262:2025), engine-version support matrix, syllabus revision, repository release tag, and external dependency manifest. Review once each academic term; do not silently replace 2025 content with newer TC39 drafts or proposals. Record new feature deltas separately as extension lessons.

\section{Final learning outcomes checklist}\label{final-learning-outcomes-checklist}

\begin{itemize}
\tightlist
\item[$\square$]
  Can locate normative language semantics by clause and trace abstract algorithms.
\item[$\square$]
  Can explain eight language types, coercion, comparison, and object invariants.
\item[$\square$]
  Can distinguish Script and Module goals, host hooks, and implementation-defined extensions.
\item[$\square$]
  Can implement tested Iterator helpers, Set algebra, Promise.try, regex escaping, and Float16 processing.
\item[$\square$]
  Can recognize attack surfaces and justify remediations with actual dataflow evidence.
\item[$\square$]
  Can build a defensible Test262-inspired traceability matrix and CI safety net.
\item[$\square$]
  Can present a capstone with reproducible tests and documented platform assumptions.
\end{itemize}

\section{Reference materials}\label{reference-materials}

\textbf{Normative source:} ECMA-262:2025, 16th edition, June 2025 (uploaded standard), clauses 1--29 and annexes as indicated. \textbf{Optional nonnormative learning infrastructure:} \url{https://github.com/tc39/test262} , \url{https://developer.mozilla.org/en-US/docs/Web/JavaScript} , \url{https://nodejs.org/api/} , \url{https://eslint.org/docs/latest/} , \url{https://codeql.github.com/docs/} . Validate hosting/runtime availability independently.

\end{document}
